% Source: https://artofproblemsolving.com/community/c3h1369019_2015_chapter_sprint_29
% Stable identity: aops-topic-1369019-q29
Write $\dfrac{1}{98}=\dfrac{1}{100}\cdot\dfrac{1}{1-\dfrac{2}{100}}=\dfrac{1}{100}\left(1+\dfrac{2}{100}+\dfrac{4}{100^2}+\dfrac{8}{100^3}+\dfrac{16}{100^4}+\cdots\right)$.

So the decimal is built from the powers of $2$ placed in successive pairs of digits:
\[0.01\,02\,04\,08\,16\,32\,64\,\ldots\]
The terms are small enough that no carries reach the first ten digits: the terms from $\dfrac{32}{100^6}$ on add up to less than $\dfrac{33}{100^6}$, which affects only the $11$th digit and beyond. (Indeed $\dfrac{1}{98}=0.0102040816326530\ldots$.) The digits are $0,1,0,2,0,4,0,8,1,6,\ldots$, so the $10$th digit is $\boxed{6}$.
